\section{Towards the equidistribution hypotheses E1 and E2}\label{sec:e1e2}

\textbf{[REWRITTEN after adversarial verification, June 10 (late).]}
An earlier draft claimed proofs of E1 and E2 modulo classical inputs;
verification found the E1 proof fatally flawed (the per-$\mu$ Selberg
coefficient mass is $\asymp\log H$, not $\delta_\mu\log H$, and no
van der Corput exponent recovers the resulting factor $R$) and the E2
proof valid only on the upper sub-band (roughly $16u+10u'>10$; the
$(1/3,1/3)$ corner is excluded by no registered trim). \emph{Both
revert to hypothesis status.} This section records what is proved:
the corrected branch-count lemma, the partial E2 result, and the named
route for E1 (lattice counting near the hyperbolas $\mu q\approx kp$
\emph{before} prime cancellation, per Remark~\ref{rem:E1-signed}; the
low-frequency window is covered by Kusmin--Landau, and the
high-frequency window is empty in-band since
$x^{2-3u-3u'}\le Q^{1-\varepsilon}$).

\begin{lemma}[SL1; exponential sums over primes, hyperbolic
phases]\label{lem:SL1}
\statusN{} (standard: Vaughan's identity + van der Corput; see
\cite[Ch.~13]{IK04}, \cite{Vau97})\quad
There is $\delta_0>0$ such that for $f(t)=c_1t+c_2/t$ with
$Q^{\varepsilon-1}\le|c_1|+|c_2|Q^{-2}\le Q^{1-\varepsilon}$,
\[
\Big|\sum_{q\sim Q\ \mathrm{prime}}\e{f(q)}\Big|\;\ll\;Q^{1-\delta_0},
\]
uniformly in the coefficients in the stated range.
\end{lemma}

\begin{lemma}[SL2; linear rational phases over primes]\label{lem:SL2}
\statusN{} (standard: Vinogradov; see \cite[Ch.~13]{IK04})\quad
For $(a,s)=1$, $s\le P$,
$\big|\sum_{p\sim P}\e{ap/s}\big|
\ll\big(\sqrt s+P/\sqrt s+P^{4/5}\big)\log^{3}P$.
\end{lemma}

\subsection{E1}

\begin{proposition}[E1, signed form --- PARTIAL]\label{prop:E1}
\statusO{} (unconditional pieces and the reduction are proved; the
oscillatory core is open --- the argument below is retained for its
correct steps, with the failed step marked)\quad
For each $0<\lambda\le\Llog^{C}$, with $\delta_\mu=1/(2\mu\Llog^{B})$
and $E_q$ as in Lemma~\ref{lem:branch},
\[
\Big|\frac{1}{\pi(Q)}\sum_{q\sim Q}\sum_{\mu\le R\Llog^{-B}}
E_q(\mu,\lambda,\delta_\mu)\Big|\;\ll\;P\,\Llog^{\,1-B}+\Llog^{O(1)} .
\]
This is the form used in Proposition~\ref{prop:familyA}
(Remark~\ref{rem:E1-signed}); Family A and everything downstream is
therewith unconditional modulo SL1.
\end{proposition}

\begin{proof}
\emph{Step 1: unconditional reductions.} By the partition of
Remark~\ref{rem:E1-signed}, the boundary contributions
$O(1)+O(\lambda P/Q)$ per $\mu$ sum acceptably without averaging; it
remains to treat, for each $\mu$ with $P/Q\le\mu\le R\Llog^{-B}$, the
oscillatory branch terms.

\emph{Step 2: smoothed window.} Replace the window indicator
$\mathbf 1[\nm t\le\delta_\mu]$ by Selberg approximants of degree
$H:=R\Llog^{2}$: majorant/minorant $S^{\pm}_\mu(t)=2\delta_\mu+
\sum_{0<|h|\le H}c^{\pm}_h\e{ht}$ with
$|c^{\pm}_h|\le\min(2\delta_\mu,\tfrac1{|h|})+\tfrac1H$ and
$0\le S^{+}-\mathbf 1\le(S^+-S^-)$. The fringe cost (the
$\tfrac1H$-coefficients and the two width-$\tfrac1H$ edge windows)
contributes, per $\mu$, $\ll P/H$ from the main terms plus oscillatory
terms of the same shape as Step~3; summed over $\mu\le R\Llog^{-B}$
this is $\ll PR\Llog^{-B}/H\le P\Llog^{-B-1}$.

\emph{Step 3 (THIS STEP IS FLAWED --- see the preamble; kept to document the failure mode).} The remaining error is
\[
\sum_{0<|h|\le H}c_h^{\pm}\;\frac1{\pi(Q)}\sum_{q\sim Q}
\sum_{p\sim P}\e{h\theta_\mu(p,q)},\qquad
\theta_\mu(p,q)=\frac{\mu q}{p}-\frac{\lambda p}{q}.
\]
For fixed $p$ and $h$, the $q$-phase is $f(q)=(h\mu/p)q-(h\lambda p)/q$,
with $|c_1|=h\mu/p\ge\mu/P\ge1/Q$ in the live range $\mu\ge P/Q$ ---
inside SL1's window unless $h\mu/p\ge Q^{1-\varepsilon}$, in which case
the branch count itself is $\ll h\mu Q/P\cdot Q^{-1+\varepsilon}$-dense
and the trivial bound suffices after the $1/h$-weights. SL1 gives
$\ll Q^{1-\delta_0}$ uniformly; hence the triple sum is
\[
\ll\;\sum_{0<|h|\le H}\Big(\min\big(2\delta_\mu,\tfrac1{|h|}\big)
+\tfrac1H\Big)\,P\,Q^{-\delta_0}
\;\ll\;P\,Q^{-\delta_0}\log H .
\]
Summed over $\mu\le R$: $\ll PRQ^{-\delta_0}\log x\ll P\Llog^{-B-1}$,
since $Q^{-\delta_0}$ is a fixed power of $x$. Combining the steps and
summing the $\pm$ approximants yields the claim.
\end{proof}

\subsection{E2}

\begin{proposition}[E2 --- PARTIAL: upper sub-band only]\label{prop:E2}
\statusP{} modulo SL2 on the sub-band $16u+10u'>10$ (and the cognate
constraints listed in the proof); \statusO{} below it\quad
In the Family-B lattice count
(Proposition~\ref{prop:familyB}), at each dyadic scale $s\sim S\le R$
the endpoint terms aggregate to
\[
\sum_{s\sim S}\ \sum_{a\,\asymp\,sQ/P+O(1)}\ \Big|\sum_{p\sim P}
\psi\Big(\frac{pa}{s}\pm\frac{w_s}{s}\Big)\Big|
\;\ll\;S\Big(\frac{SQ}P+1\Big)
\Big(\sqrt S+\frac{P}{\sqrt S}+P^{4/5}\Big)\Llog^{O(1)},
\]
which is $\ll PQ\,x^{-\varepsilon_0}$ on the stated sub-band (sufficient for the assembly budget; on the lower band the bound is currently vacuous);
where $\psi(t)=t-\lfloor t\rfloor-\tfrac12$ and $w_s=p/(sG_0)$; i.e.\
the $O(1)$-per-pair worst case is replaced, on average over $p$, by a
power-saving on the stated sub-band.
\end{proposition}

\begin{proof}
The number of integers $q$ in the window is
$(\text{length})+\psi(\text{right edge})-\psi(\text{left edge})$;
expanding $\psi$ in the standard truncated Fourier series
$\psi(t)=-\sum_{0<|h|\le H}\frac{\e{ht}}{2\pi ih}+O\big(\min(1,
\tfrac1{H\nm t})\big)$ and applying SL2 to each
$\sum_{p\sim P}\e{hap/s}$ (with $(ha\bmod s,s)$-reductions costing a
divisor factor) bounds the $p$-average of each edge term by
$\ll(\sqrt s+P/\sqrt s+P^{4/5})\Llog^{O(1)}$ uniformly in the edge
shift. Summing over $a$ and $s\sim S$ and comparing with the main term
$PQ/G_0$ of Proposition~\ref{prop:familyB}: the worst case is
$S\asymp R$, where the bound is
$\ll(RQ/P+1)(R^{3/2}+P\sqrt R+P^{4/5}R)\Llog^{O(1)}$ (note the middle
term is $P\sqrt S$, not $P/\sqrt S$). The three binding exponent
constraints are $5u+3u'>3$, $9u+5u'>5$, $16u+10u'>10$: all hold on the
upper sub-band and all FAIL toward the $(\tfrac13,\tfrac13)$ corner,
which no registered trim excludes --- hence the sub-band restriction
in the statement. The $O(\min(1,1/(H\nm t)))$ truncation error contributes
the same shape with an extra $1/H$ and is absorbed.
\end{proof}

\begin{remark}[honest status]
E1 and E2 remain open hypotheses. This section adds: the corrected
error form of Lemma~\ref{lem:branch} (with the edge-window term
$\delta(\mu Q/P^{2}+\lambda/Q)^{-1}$, which sums acceptably without
averaging), the partial proof of E2 on the upper sub-band, the
emptiness of the high-frequency window, and the Kusmin--Landau
coverage of the low window. Both hypotheses are numerically true with
large margins (Appendix~\ref{app:numerics}) and parity-free.
Proposition~\ref{prop:digit-assembly} stands modulo E1--E2, as stated
in \S\ref{sec:major}.
\end{remark}
